% !TEX root = ../main.tex
\chapter{Tracer and TraceIR}
\label{ch:tracer}

The tracer is the front end of \Lucid's compile path. Given a
model and a representative input shape, it runs the model's
forward pass in a special \emph{tracing mode} that records every
primitive op call as a node in an intermediate representation
(the trace IR), without executing the actual kernels. The
output is a graph that the MpsBuilder
(\chref{ch:mps_builder}) translates into a compiled
\MPSGraph\ executable.

This chapter describes the tracer's design, the TraceIR data
structures, the tracing protocol, the limitations of dynamic
tracing, and the diagnostic tools that expose the trace to the
user.

\section{Why Tracing}
\label{sec:tracer:why_tracing}

A compiler needs to know what the model computes before it can
compile. Two general strategies:

\begin{itemize}
  \item \textbf{Static analysis} (AOT, ``ahead of time''): parse
    the model's source code to extract the computation graph.
    This is what TorchScript attempted; it requires the user to
    write models in a restricted Python subset.
  \item \textbf{Tracing} (JIT, ``just in time''): run the model
    once on a representative input, recording every operation as
    it executes. This is what \refframework's \code{torch.compile},
    JAX's \code{jit}, and \Lucid's \code{compile} do.
\end{itemize}

Tracing is more flexible: the user writes ordinary Python and
the tracer adapts. The cost is that data-dependent control flow
(Python \code{if} based on tensor values) is invisible to the
tracer; it sees only the branch taken on the tracing input.

\section{The Trace IR}
\label{sec:tracer:trace_ir}

A \emph{trace} is a directed acyclic graph of op invocations.
Each node is an instance of \code{TraceNode}, defined in
\code{lucid/\_C/compile/TraceIR.h}:

\begin{lstlisting}[style=lucidcpp,language=C++]
struct TraceNode {
  std::string             op_name;     // canonical op name
  std::vector<size_t>     input_ids;   // upstream node ids
  std::vector<TraceTensor> outputs;    // shape+dtype only
  std::vector<Attribute>  attrs;       // op-specific scalars
  uint32_t                version;     // schema version
};

struct TraceTensor {
  Shape   shape;
  Dtype   dtype;
  Device  device;
  // no actual data --- only metadata
};

struct Trace {
  std::vector<TraceNode>  nodes;       // topologically ordered
  std::vector<size_t>     input_ids;   // graph inputs
  std::vector<size_t>     output_ids;  // graph outputs
  std::vector<Attribute>  metadata;    // trace-level info
};
\end{lstlisting}

The IR is dataflow: nodes are operations, edges (via
\code{input\_ids}) are tensor flow. The IR is topologically
ordered so that emission can walk it linearly.

\subsection{Shape-only tensors}
\label{ssec:tracer:shape_only}

A \code{TraceTensor} carries metadata (shape, dtype, device) but
no actual data. The tracer never materialises kernels; it only
tracks what the kernels would produce.

This is what makes tracing cheap. The trace pass for a 100M-
parameter Transformer takes tens of milliseconds, regardless of
the eventual compilation cost.

\subsection{Op attributes}
\label{ssec:tracer:attributes}

Ops with non-tensor arguments (e.g., \code{conv2d}'s
\code{stride}, \code{padding}) store them as
\code{Attribute}s: scalar values or small sequences attached to
the node. The emitter passes these through to the corresponding
\MPSGraph\ call.

\section{The Tracing Protocol}
\label{sec:tracer:protocol}

The tracer hooks into the dispatcher. When tracing is active:

\begin{enumerate}
  \item Each \code{Dispatcher::unary} / \code{binary} / etc.\ call
    is intercepted.
  \item Instead of executing the op, the tracer creates a new
    \code{TraceNode} with the op name, input node ids, and
    appropriate attributes.
  \item The output is a fresh \code{TraceTensor} with the shape
    and dtype that the op would produce (computed via the op's
    shape function, not via execution).
  \item The output is returned to the caller as a special
    \code{TracedTensor} that wraps the trace node id.
\end{enumerate}

The \code{TracedTensor} flows through the model's forward like
an ordinary tensor; every operation on it produces a new
\code{TracedTensor} via the same hook.

\subsection{Shape inference}
\label{ssec:tracer:shape_inference}

Each op's schema (\chref{ch:op_registry}) includes a shape
function: given input shapes, return output shape. The tracer
calls it to determine the output \code{TraceTensor}. Shape
functions are typically a few lines (most ops preserve or
broadcast input shapes) and are validated by the parity tests.

\subsection{Entering and exiting trace mode}
\label{ssec:tracer:enter_exit}

\code{compile()} enters trace mode via a context manager:

\begin{lstlisting}[style=lucidcpp,language=C++]
{
  TraceContext ctx;
  TracedTensor input_traced = make_traced(input_shape, input_dtype);
  // Run the user's forward with input replaced by traced tensor
  TracedTensor output_traced = model.forward(input_traced);
  // Extract the recorded trace
  Trace trace = ctx.extract();
}
\end{lstlisting}

Outside the context, the dispatcher's hooks are no-ops; normal
eager execution proceeds.

\section{Dynamic Control Flow}
\label{sec:tracer:dynamic_control_flow}

The fundamental limitation of tracing: Python \code{if}/\code{while}
based on tensor \emph{values} is invisible. The tracer sees only
the branch taken on the tracing input.

\subsection{Example: a data-dependent branch}
\label{ssec:tracer:branch_example}

\begin{lstlisting}[style=lucid,language=LucidPython]
def forward(x):
    if x.mean() > 0:
        return x.relu()
    else:
        return -x.relu()
\end{lstlisting}

The tracer sees one branch on its first call. A future call where
the branch flips will produce wrong results unless the cache
recompiles for that input.

The user-side mitigation: rewrite the branch in tensor-level form
using \code{where}:

\begin{lstlisting}[style=lucid,language=LucidPython]
def forward(x):
    pos = x.mean() > 0
    return lucid.where(pos, x.relu(), -x.relu())
\end{lstlisting}

\code{where} is a tensor op that the tracer captures correctly.
The user must compute both branches (the cost is the union of
both branch costs), but the trace is consistent.

\subsection{Example: a loop with variable count}
\label{ssec:tracer:loop_example}

A loop whose iteration count depends on tensor values (e.g., a
training-time iterative algorithm that converges when residual
is small) is also opaque. The tracer sees the count from the
tracing input.

The user-side mitigation: trace the inner step, run the loop in
Python; each iteration calls the compiled step. The compile
boundary is at the inner step, not the loop.

\section{Diagnostic Tools}
\label{sec:tracer:diagnostics}

The tracer exposes its output for diagnostics.

\subsection{dump\_trace}
\label{ssec:tracer:dump_trace}

\code{lucid.compile.dump\_trace(compiled\_module)} prints the
recorded trace IR in a human-readable form:

\begin{verbatim}
%0 = input shape=(8, 1024, 768) dtype=float32
%1 = layer_norm(%0, normalized_shape=(768,)) -> (8, 1024, 768)
%2 = matmul(%1, weight=(768, 2304)) -> (8, 1024, 2304)
%3 = split(%2, dim=-1, sections=3) -> [(8,1024,768), ...]
%4 = reshape(%3[0], (8, 1024, 12, 64)) -> (8, 1024, 12, 64)
...
\end{verbatim}

Each line shows the SSA-style assignment, the op name, the
inputs, and the output shape. The user can verify that the trace
matches their model's intent before investigating performance.

\subsection{Trace size}
\label{ssec:tracer:trace_size}

Typical trace sizes:

\begin{itemize}
  \item GPT-2-base forward + backward: $\sim 2400$ nodes.
  \item ViT-B/16 forward + backward: $\sim 800$ nodes.
  \item ResNet-50 forward + backward: $\sim 1200$ nodes.
\end{itemize}

The trace is large but the per-node cost is small (under a
microsecond to record), so tracing time is dominated by the
shape-inference work, not the recording.

\section{Tracing Edge Cases}
\label{sec:tracer:edge_cases}

A short catalogue of the cases that have caused user confusion.

\subsection{Python scalar arguments}
\label{ssec:tracer:py_scalar}

Op arguments that are Python scalars (\code{int}, \code{float},
\code{bool}) are baked into the trace as constants. A
\code{linear(x, 0.5)} traces with the 0.5 hard-coded; calling
the compiled module with a different scalar would not change the
behaviour because the trace does not see the new value.

The mitigation: pass scalars as 0-d tensors when they need to be
trace-time-variable.

\subsection{Lists of tensors}
\label{ssec:tracer:tensor_lists}

A function returning \code{[t1, t2, t3]} traces as three
distinct outputs. The trace knows the list structure but the
\MPSGraph\ executable returns a flat list; the
\code{CompiledModule} rewraps in the original structure.

\subsection{Nested modules}
\label{ssec:tracer:nested_modules}

The tracer traces through nested \code{Module} calls naturally:
\code{block(x)} eventually calls primitive ops, and those are
the nodes that appear in the trace. Module identity is lost; the
trace contains primitive ops only.

\subsection{Print statements}
\label{ssec:tracer:print_statements}

A \code{print(tensor)} inside the forward forces a value read,
which the tracer sees as a request to materialise the
\code{TracedTensor}. Materialisation is impossible during
tracing (the tensor has no value), so the trace raises.

The user-side: remove print statements from the forward; use
hooks or the post-compile path for inspection.

\section{Re-tracing}
\label{sec:tracer:retracing}

A new input shape (or, more generally, anything that changes the
trace's structure) requires re-tracing. The cost is:

\begin{itemize}
  \item Tracing itself: tens of milliseconds.
  \item MpsBuilder + \MPSGraph\ compile: hundreds of
    milliseconds.
\end{itemize}

For shape-stable training loops, re-tracing happens once at
startup. For shape-variable workloads (variable-length sequences
without padding), re-tracing can happen per batch, defeating
the compile benefit. The standard mitigation is to pad to fixed
shape.

\section{Implementation Notes}
\label{sec:tracer:impl_notes}

The tracer lives in \code{lucid/\_C/compile/Tracer.\{h,cpp\}}.
The IR types are in \code{TraceIR.h}. The dispatcher hook is
installed via a thread-local \code{TraceContext} pointer that
the dispatcher checks at every call.

\subsection{Performance overhead}
\label{ssec:tracer:perf_overhead}

The per-op overhead during tracing is roughly $5\,\mu s$
(checking the context, computing the output shape, allocating a
\code{TraceNode}). For a typical model with a few thousand ops,
the tracing pass is in the tens-of-milliseconds range.

This is small relative to the subsequent \MPSGraph\ compilation
(hundreds of milliseconds to seconds), so tracing is rarely the
bottleneck of cold start.

\subsection{Memory footprint}
\label{ssec:tracer:mem_footprint}

The trace IR itself takes a few hundred kilobytes for a typical
model. The trace is held by the cache for the lifetime of the
compiled module; for users with very large or many compiled
modules, this is observable in process memory.

\section{Summary and Forward References}
\label{sec:tracer:summary}

The tracer runs the user's forward in shape-only mode and
records each primitive op as a \code{TraceNode} in the
\code{TraceIR}. The trace is the input to the MpsBuilder
(\chref{ch:mps_builder}), which translates it to a compiled
\MPSGraph\ executable. The tracing approach is more flexible
than static analysis but cannot handle data-dependent control
flow; the user-side fix is to rewrite branches using tensor ops
like \code{where}.

The next chapter, \chref{ch:mps_builder}, treats the MpsBuilder
that consumes the trace and emits the \MPSGraph.
\chref{ch:op_emitters} catalogues the per-op emitters.

\begin{lucidnote}
For the user: the tracer is invisible at the API surface.
\code{compile()} handles the tracing internally; the user just
sees a CompiledModule. The diagnostic dumps (\code{dump\_trace})
are available for the rare cases where the user needs to inspect
what the compiler saw.
\end{lucidnote}
